\section{How the content is divided amongst the chapters}

Broadly, the thesis divides into three parts. The first is groundwork: the
stochastic processes and the methods used to solve them, together with one
constant that took long enough to pin down to deserve a chapter of its own.
The second is the birth--death--catastrophe process --- the single infected
cell, the population of cells built from it, a two-type version, and what the
whole construction says about bursting pathogenesis. The third lets the rates
move, and is where the mathematics of the earlier chapters gives out. Including
this introduction and the conclusion, these parts are divided amongst ten
chapters. Below is a summary of the content and message of each.

\subsection*{Chapter 2: Mathematical introduction}

Much of the work presented in this thesis surrounds topics of stochastic
process theory, and specifically branching processes. This chapter assembles
the objects and the methods: Markov chains in discrete and continuous time, the
branching processes that matter here, and the generating-function machinery ---
forward and backward equations, means and higher moments, hitting and
extinction probabilities, limiting conditional means, and the method of
characteristics --- each developed once, on whichever process shows it most
plainly. Two extensions then leave the time-homogeneous setting, in
anticipation of \mextref{var:ch:rates}{Chapter~9}: processes whose rates vary
with time, and coupled systems in which a differential equation and a Markov
chain each drive the other.

Two of its appendices are more than review. The first records how heavily
survival near criticality depends on the size of the founding cohort
(\mextref{m:app:smallpops}{Section~2.A}), which is the fact that almost
everything downstream leans on. The second solves a family of absorption
models, in which particles are drawn into a compartment and may then replicate
inside it; the hardest member of that family is solved here for the first time
(\mextref{m:app:absorption}{Section~2.B}).

\subsection*{Chapter 3: Galton--Watson branching and the constant
\texorpdfstring{$A(p)$}{A(p)}}

For the subcritical binary Galton--Watson process the survival probability
decays as $S_n\sim A(p)(2p)^n$, and that single constant also fixes the
population conditioned on survival, whose mean converges to $1/A(p)$. The
constant exists because the limit can be rewritten as a convergent infinite
product. Rescaling the survival recursion to the logistic map identifies it
with a value of the associated Koenigs linearising function, and a theorem of
Becker and Bergweiler shows that function to be hypertranscendental
\cite{BeckerBergweiler1995}. That does not settle the closed-form question for
$p\mapsto A(p)$, but it does explain why a long search for a formula came to
nothing. What remains is practical: the product, together with the
near-critical asymptotic $A(p)\sim2(1-2p)$, describes the constant across the
whole subcritical range.

This is the chapter furthest from plague, and it is kept because the question
it answers is a real one and the answer is complete.

\subsection*{Chapter 4: The birth--death--catastrophe process}

Around the 1980s, work was being done into growth models that incorporated
stochastic declines. Out of this came the birth--death--catastrophe model as we
know it
\cite{brockwell1982birth,brockwell1985extinction,cairns2004extinction,dori2018family}:
an extension of the standard birth--death process in which, from every state
apart from $0$, a catastrophe event can take place. We were interested in this
model for its potential application to the development of an infection within
cells that burst.

This chapter defines the process and derives the single-cell theory. The
calculation closes because the probability of no rupture yet, $I(t)$, obeys an
autonomous Riccati equation; the two roots $a$ and $b$ of its characteristic
quadratic then carry every closed form in the thesis. Cell-fate probabilities,
the mean and second moment of the intracellular load, the mean release and the
first two conditional moments of the burst size all follow in that one
coordinate. Some of the results have been found by others before, but several
--- the second moment expressed through $I$, the variances, and the joint
load--release moments --- go beyond what the existing literature had recorded.

\subsection*{Chapter 5: Distribution theory, quasi-stationarity and burst
statistics}

The previous chapter gives the moments; this one gives the laws. The
intracellular load is geometric at every time, with a ratio that slides towards
$1/a$; the quasi-stationary law --- the behaviour of the process count given
that fixation has not yet occurred --- is that endpoint; and the burst size,
averaged over the instant of rupture, is the same law again.

That the two coincide is the most satisfying result in the thesis, because the
two constructions share no step: one conditions a surviving trajectory on never
being absorbed, the other averages over rupture times, and the identity holds
because size-biasing telescopes the occupation integral onto the endpoint of
the process's own geometric ratio. Two probes then mark the boundary of the
structure. Several founders sharing a cell leave it intact; chaining one cell's
burst into the next destroys it as soon as internal extinction is possible, and
that case remains open.

\subsection*{Chapter 6: Burst-aware within-host dynamics}

This chapter embeds the cell in a population. The survival and release
statistics of the previous two chapters become kernels indexed by cell age, and
the resulting renewal system is exact in the mean rather than an approximation
to it. Three consequences follow. The standard model of within-host virus
dynamics is a projection of that system, valid one exponential phase at a time,
so a fitted release rate is a property of the phase and not of the organism,
and a model that already carries a hand-fitted eclipse alongside a bursting
pathogen is fitting the same delay twice. No fit of the two population
constants recovers the three intracellular rates, so a population fit is worth
doing only alongside a single-cell measurement. And at fixed total output under
linear clearance, the release schedule is invisible to $R_0$ while remaining
visible to establishment: bursting strictly lowers the extinction probability
of a small inoculum if and only if the flooding parameter $L=a(1-b)$ exceeds
one.

\subsection*{Chapter 7: Type-specific catastrophe rates}

In 2011 Antal and Krapivsky \cite{antal2011exact} published an exact solution
to the two-type continuous time birth--death process with one-way mutation. We
were able to use their method, with some changes of variable, to obtain the
equivalent result once catastrophes are allowed.

Type~1 converts irreversibly into type~2, either type may trigger a single
global catastrophe, and because the catastrophe rate depends on the whole path,
the natural object is an occupation-time transform rather than a
terminal-count generating function. The backward system is triangular, and
after linearisation the surviving Riccati equation becomes the Gauss
hypergeometric equation, giving a closed formula for the finite-time
probability of avoiding catastrophe which holds for any ordering of the
type-specific rates. Read biologically, the two types are the early and adapted
intracellular phenotypes of \yp{} described above, and the catastrophe is
irreversible loss of macrophage containment.

\subsection*{Chapter 8: Bursting and budding pathogenesis}

Work has been done by others
\cite{komarova2007viral,yuan2011stochastic,pearson2011stochastic,capistran2018extracellular}
comparing the dynamical differences between budding and bursting pathogenesis.
This chapter asks what a burst is worth when the defence facing it can be used
up --- when there is, as it were, a safety in numbers.

Inside the cell, a replicator--suppressor process treats the toxin-containing
granules of a macrophage or neutrophil as a store which is spent and never
replenished. A larger incident cohort is then disproportionately safer, because
the cohort itself depletes the store. Outside it, in the extracellular medium,
the pathogen meets complement proteins and killer cells which are likewise
finite on the relevant timescale; modelling this as a fixed removal budget
$\vartheta$ per delivery gives the surviving release in closed form,
$\mathcal{R}(\vartheta)=V_\infty a^{-\vartheta}$, so that every unit of removal
costs the burst one factor of $a$. With it comes a reversal: against a
threshold rather than a rate, the flooding criterion of the previous chapters
flips, and bursting beats matched-mean budding precisely when $L<1$. The reason
is a single fact about variance --- extinction probability rewards low
offspring variance, whereas the surplus above a threshold is convex and rewards
high variance --- so the sign of $L-1$ drives both criteria, and drives them
opposite ways.

\subsection*{Chapter 9: Models with non-constant rates}

Every model carried through to a result in the preceding chapters holds its
rates fixed, and that assumption is what makes the generating functions close.
It is also the first thing an ecologist will question. This chapter asks what
happens when a rate is allowed to move, and it came from a different appetite
from the rest of the thesis; it marks plainly which of its parts are results
and which are specifications.

In the spirit of Yule's paper and of more recent work with birth--death
processes \cite{nee1994extinction}, the first model gives a genus an
inhomogeneous origination rate, driven by the room a population has left to
expand into under logistic growth. The species count stays geometric and
depends on the rate law only through its integral; under the logistic law that
integral converges, leaving a proper limit law rather than merely a finite
late-time mean. The second treats Pimentel's fly experiment
\cite{pimentel1965selection,hutchingson2018ecological}, in which two species
were shown to coexist in one niche by oscillation, because a population pushed
towards extinction evolves in response to interspecific pressure more than
intraspecific; here the state at which the chain is equally likely to step up
as down is available in closed form, and it is unstable, which is a different
thing from a process approaching an equilibrium. Two further models, in which a
population accumulates a potential and then loses it in a catastrophe, are
carried only as far as simulation.

The chapter's general observation is worth repeating, because it is not the
obvious one: what costs the generating function is not feedback as such but
continuous bending of a rate. A process may act on its own circumstances as
much as it likes and stay solvable, provided what survives the interaction is a
random time or a random count rather than a moving rate.

\subsection*{Chapter 10: Conclusion}

The conclusion is a reckoning rather than a summary. It sets out what each
chapter established and what it left standing; returns to the four questions posed
above and answers them as far as the work allows;
gives the quadrant argument, a small piece of reasoning about why an organism
might want to be eaten early and not late, which never found a home in a
chapter of its own; and says plainly what the thesis does not establish,
together with the four open problems worth taking first.
